\section{Existence and Uniqueness of Solutions}

In this section, we establish the well-posedness of the stochastic advection--diffusion equation. 
The proof relies on standard techniques from the theory of stochastic evolution equations, 
including Galerkin approximations, energy estimates, and Grönwall's inequality.

\subsection{Assumptions}

Throughout this section, we assume that:

\begin{itemize}
\item the domain \(D\subset\mathbb{R}^d\) is bounded with Lipschitz boundary;
\item the diffusion coefficient satisfies
\[
\kappa>0;
\]
\item the velocity field satisfies
\[
\mathbf b\in L^\infty(D)^d;
\]
\item the initial condition satisfies
\[
y_0\in L^2(\Omega;H);
\]
\item the control satisfies
\[
u\in L^2\bigl(\Omega;L^2(0,T;H)\bigr).
\]
\end{itemize}

Under these assumptions, the bilinear form \(a(\cdot,\cdot)\) defined by \eqref{eq:bilinear_form} 
is continuous on \(V\times V\).

\begin{lemma}
There exists a constant \(C_a>0\) such that
\[
|a(y,v)|
\le
C_a
|y|_V
|v|_V,
\qquad
\forall y,v\in V.
\]
\end{lemma}

\begin{proof}
The result follows from the Cauchy--Schwarz inequality and the boundedness of the velocity field 
\(\mathbf b\).
\end{proof}

Moreover, the bilinear form satisfies the following G{\aa}rding inequality.

\begin{lemma}[G{\aa}rding inequality]
There exists a constant \(C_G>0\) such that
\[
a(v,v)
\ge
\kappa
|\nabla v|_{L^2(D)}^2
C_G
|v|_{L^2(D)}^2,
\qquad
\forall v\in V.
\]
\end{lemma}

\begin{proof}
Using the definition of \(a(\cdot,\cdot)\), we obtain
\[
a(v,v)
=
\kappa
\int_D |\nabla v|^2,dx
+
\int_D (\mathbf b\cdot\nabla v),v,dx.
\]

Applying the Cauchy--Schwarz and Young inequalities yields
\[
\left|
\int_D (\mathbf b\cdot\nabla v),v,dx
\right|
\le
C
|v|*{L^2(D)}
|\nabla v|*{L^2(D)}
\]

and therefore

\[
a(v,v)
\ge
\kappa
|\nabla v|_{L^2(D)}^2
-
C_G
|v|_{L^2(D)}^2.
\]

This concludes the proof.
\end{proof}

\subsection{Existence and Uniqueness}

We are now ready to state the main well-posedness result.

\begin{theorem}
Under the assumptions stated above, there exists a unique weak solution \(y\) of problem 
\eqref{eq:state}--\eqref{eq:ic} satisfying

\[
y
\in
L^2
\Bigl(
\Omega;
L^2(0,T;V)
\Bigr)
\cap
L^2
\Bigl(
\Omega;
C([0,T];H)
\Bigr).
\]

Moreover, there exists a constant \(C>0\) independent of \(y\) such that

\[
\E
\Bigl[
\sup_{0\le t\le T}
|y(t)|_H^2
\Bigr]
+
\E
\Bigl[
\int_0^T
|y(t)|_V^2,dt
\Bigr]
\le
C
\left(
\E|y_0|_H^2
+
\E
\int_0^T
|u(t)|_H^2,dt
+
\sigma^2 T
\right).
\]
\label{thm:existence}
\end{theorem}

\begin{proof}
The proof follows the classical Galerkin method.

Let \({\varphi_i}_{i\ge1}\) be an orthonormal basis of \(H\) consisting of eigenfunctions of the 
Laplace operator. For each \(N\), we seek an approximation

\[
y_N(t)
=
\sum_{i=1}^{N}
g_i(t)\varphi_i.
\]

Substituting this expansion into the weak formulation \eqref{eq:weak_form} yields a 
finite-dimensional stochastic differential system.

Applying It\^o's formula to

\[
|y_N(t)|_H^2
\]

and using the G{\aa}rding inequality gives an energy estimate of the form

\[
\E
\Bigl[
\sup_{0\le t\le T}
|y_N(t)|_H^2
\Bigr]
+
\E
\Bigl[
\int_0^T
|y_N(t)|_V^2,dt
\Bigr]
\le C.
\]

The estimate is uniform with respect to \(N\). Consequently, weak compactness arguments 
allow the extraction of a subsequence converging to a limit \(y\).

Passing to the limit in the weak formulation shows that \(y\) is a weak solution.

To prove uniqueness, let \(y_1\) and \(y_2\) be two weak solutions and define

\[
z=y_1-y_2.
\]

Then \(z\) satisfies the homogeneous equation. Applying It\^o's formula to \(|z(t)|_H^2\) 
and using Grönwall's inequality yields

\[
z=0.
\]

Hence \(y_1=y_2\), which proves uniqueness.
\end{proof}

Theorem \ref{thm:existence} guarantees that the stochastic state equation is well posed. This 
result constitutes the foundation for the optimal control problem introduced in the next section.
